% GENERATED FILE. Edit canonical lessons, then run npm run generate:books.
% canonical-source-hash: fnv1a64:3d7b244dd083f749
% canonical-lessons: TA-C88-hospital, TA-C88-vomiting, TA-C88-headache, TA-C88-health, TA-C88-sweat

\chapter{At the Hospital}
\label{ch:ta-at-the-hospital}
\hlchaptermodality{88}

\begin{chapteropening}
\textbf{By the end of this chapter:} \emph{I can ask for the hospital, say I feel sick or have a headache, and talk about health and sweat.}

Complete the last lesson of chapter 88: i can ask for the hospital, say I feel sick or have a headache, and talk about health and sweat.
\end{chapteropening}

\section[\texorpdfstring{\ta{மருத்துவமனை} (maruttuvamaṉai)}{மருத்துவமனை (maruttuvamaṉai)}]{\ta{மருத்துவமனை} (maruttuvamaṉai) — a hospital}

\label{lesson:TA-C88-hospital}



\begin{quote}
Before the new one: say the Tamil for weariness, then the Tamil for rest.
\end{quote}

\subsection*{You'll want to know: \ta{மருத்துவமனை}}
\textbf{\ta{மருத்துவமனை}} (\emph{maruttuvamaṉai}) — \textquotedblleft{}a hospital\textquotedblright{}, literally \textquotedblleft{}medicine-house\textquotedblright{}.

\textbf{\ta{மனை}} is \textquotedblleft{}a house, a dwelling\textquotedblright{}. \textbf{\ta{மருத்துவமனை} \ta{எங்கே}?} — \textquotedblleft{}where is the hospital?\textquotedblright{}.

\textbf{In the family, and next door}

\noindent
\begin{tabularx}{\linewidth}{@{}>{\raggedright\arraybackslash}X>{\raggedright\arraybackslash}X@{}}
\toprule
\textbf{} & \textbf{word} \\
\midrule
Kannada & \ka{ಆಸ್ಪತ್ರೆ} (\emph{āspatre}) \\
Telugu & \te{ఆసుపత్రి} (\emph{āsupatri}) \\
Malayalam & \ml{ആശുപത്രി} (\emph{āśupatri}) \\
Hindi (neighbour) & \dv{अस्पताल} (\emph{aspatāl}) \\
English & a hospital \\
\bottomrule
\end{tabularx}

\begin{grammarlens}[title={What you've built}]
A place to ask the way to.
\end{grammarlens}

\subsection*{Your turn}
\begin{itemize}
\raggedright
  \item \emph{Say it:} \emph{maruttuvamaṉai}
  \item \emph{Say it:} \emph{maruttuvamaṉai}, once more
  \item \emph{Say it:} ask \emph{maruttuvamaṉai eṅkē?}
  \item \emph{Recall:} say the Tamil for weariness, then the Tamil for rest, then say \emph{maruttuvamaṉai} again
\end{itemize}

\begin{tcolorbox}[breakable,colback=teal!4,colframe=teal!35!black,title={Before you move on}]
What does \ta{மருத்துவமனை} mean? (\textquotedblleft{}A hospital\textquotedblright{}.) Say it once more.
\end{tcolorbox}
\section[\texorpdfstring{\ta{வாந்தி} (vānti)}{வாந்தி (vānti)}]{\ta{வாந்தி} (vānti) — vomiting, nausea}

\label{lesson:TA-C88-vomiting}



\begin{quote}
Before the new one: say the Tamil for rest, then the Tamil for a hospital.
\end{quote}

\subsection*{You'll want to know: \ta{வாந்தி}}
\textbf{\ta{வாந்தி}} (\emph{vānti}) — \textquotedblleft{}vomiting\textquotedblright{}.

\textbf{\ta{வாந்தி} \ta{வருகிறது}} — \textquotedblleft{}I feel sick\textquotedblright{}, literally \textquotedblleft{}vomiting is coming\textquotedblright{}.

\textbf{In the family, and next door}

\noindent
\begin{tabularx}{\linewidth}{@{}>{\raggedright\arraybackslash}X>{\raggedright\arraybackslash}X>{\raggedright\arraybackslash}X@{}}
\toprule
\textbf{} & \textbf{word} & \textbf{same root} \\
\midrule
Kannada & \ka{ವಾಂತಿ} (\emph{vānti}) & yes \\
Malayalam & \ml{ഛർദ്ദി} (\emph{chardi}) &  \\
English & vomiting, nausea &  \\
\bottomrule
\end{tabularx}

\begin{grammarlens}[title={What you've built}]
A symptom a doctor will ask about.
\end{grammarlens}

\subsection*{Your turn}
\begin{itemize}
\raggedright
  \item \emph{Say it:} \emph{vānti}
  \item \emph{Say it:} \emph{vānti}, once more
  \item \emph{Say it:} say \emph{vānti varukiṟatu}
  \item \emph{Recall:} say the Tamil for rest, then the Tamil for a hospital, then say \emph{vānti} again
\end{itemize}

\begin{tcolorbox}[breakable,colback=teal!4,colframe=teal!35!black,title={Before you move on}]
What does \ta{வாந்தி} mean? (\textquotedblleft{}Vomiting, nausea\textquotedblright{}.) Say it once more.
\end{tcolorbox}
\section[\texorpdfstring{\ta{தலைவலி} (talaivali)}{தலைவலி (talaivali)}]{\ta{தலைவலி} (talaivali) — a headache}

\label{lesson:TA-C88-headache}



\begin{quote}
Before the new one: say the Tamil for a hospital, then the Tamil for vomiting.
\end{quote}

\subsection*{You'll want to know: \ta{தலைவலி}}
\textbf{\ta{தலைவலி}} (\emph{talaivali}) — \textquotedblleft{}a headache\textquotedblright{}.

It is \textbf{\ta{தலை}}, \textquotedblleft{}head\textquotedblright{}, and \textbf{\ta{வலி}}, \textquotedblleft{}pain\textquotedblright{}, written as one word. \textbf{\ta{எனக்குத்} \ta{தலைவலி}} — \textquotedblleft{}I have a headache\textquotedblright{}.

\textbf{In the family, and next door}

\noindent
\begin{tabularx}{\linewidth}{@{}>{\raggedright\arraybackslash}X>{\raggedright\arraybackslash}X@{}}
\toprule
\textbf{} & \textbf{word} \\
\midrule
Kannada & \ka{ತಲೆನೋವು} (\emph{talenōvu}) \\
Malayalam & \ml{തലവേദന} (\emph{talavēdana}) \\
Hindi (neighbour) & \dv{सिरदर्द} (\emph{sirdard}) \\
English & a headache \\
\bottomrule
\end{tabularx}

\begin{grammarlens}[title={What you've built}]
Head and pain, joined.
\end{grammarlens}

\subsection*{Your turn}
\begin{itemize}
\raggedright
  \item \emph{Say it:} \emph{talaivali}
  \item \emph{Say it:} \emph{talaivali}, once more
  \item \emph{Say it:} say \emph{eṉakkut talaivali}
  \item \emph{Recall:} say the Tamil for a hospital, then the Tamil for vomiting, then say \emph{talaivali} again
\end{itemize}

\begin{tcolorbox}[breakable,colback=teal!4,colframe=teal!35!black,title={Before you move on}]
What does \ta{தலைவலி} mean? (\textquotedblleft{}A headache\textquotedblright{}.) Say it once more.
\end{tcolorbox}
\section[\texorpdfstring{\ta{உடல்நலம்} (uṭalnalam)}{உடல்நலம் (uṭalnalam)}]{\ta{உடல்நலம்} (uṭalnalam) — health}

\label{lesson:TA-C88-health}



\begin{quote}
Before the new one: say the Tamil for vomiting, then the Tamil for a headache.
\end{quote}

\subsection*{You'll want to know: \ta{உடல்நலம்}}
\textbf{\ta{உடல்நலம்}} (\emph{uṭalnalam}) — \textquotedblleft{}health\textquotedblright{}, literally \textquotedblleft{}body-wellbeing\textquotedblright{}.

\textbf{\ta{நலம்}} is the wellbeing of your first greetings; \textbf{\ta{உடல்}} is \textquotedblleft{}the body\textquotedblright{}.

\textbf{In the family, and next door}

\noindent
\begin{tabularx}{\linewidth}{@{}>{\raggedright\arraybackslash}X>{\raggedright\arraybackslash}X@{}}
\toprule
\textbf{} & \textbf{word} \\
\midrule
Kannada & \ka{ಆರೋಗ್ಯ} (\emph{ārōgya}) \\
Telugu & \te{ఆరోగ్యం} (\emph{ārōgyaṁ}) \\
Hindi (neighbour) & \dv{सेहत} (\emph{sehat}) \\
English & health \\
\bottomrule
\end{tabularx}

\begin{grammarlens}[title={What you've built}]
Health, built from a word you already use.
\end{grammarlens}

\subsection*{Your turn}
\begin{itemize}
\raggedright
  \item \emph{Say it:} \emph{uṭalnalam}
  \item \emph{Say it:} \emph{uṭalnalam}, once more
  \item \emph{Say it:} say \emph{uṭalnalam}
  \item \emph{Recall:} say the Tamil for vomiting, then the Tamil for a headache, then say \emph{uṭalnalam} again
\end{itemize}

\begin{tcolorbox}[breakable,colback=teal!4,colframe=teal!35!black,title={Before you move on}]
What does \ta{உடல்நலம்} mean? (\textquotedblleft{}Health\textquotedblright{}.) Say it once more.
\end{tcolorbox}
\section[\texorpdfstring{\ta{வியர்வை} (viyarvai)}{வியர்வை (viyarvai)}]{\ta{வியர்வை} (viyarvai) — sweat}

\label{lesson:TA-C88-sweat}



\begin{quote}
Before the new one: say the Tamil for a headache, then the Tamil for health.
\end{quote}

\subsection*{You'll want to know: \ta{வியர்வை}}
\textbf{\ta{வியர்வை}} (\emph{viyarvai}) — \textquotedblleft{}sweat\textquotedblright{}.

\textbf{\ta{வியர்வை} \ta{வருகிறது}} — \textquotedblleft{}I am sweating\textquotedblright{}, literally \textquotedblleft{}sweat is coming\textquotedblright{}, like \textbf{\ta{வாந்தி}}.

\textbf{In the family, and next door}

\noindent
\begin{tabularx}{\linewidth}{@{}>{\raggedright\arraybackslash}X>{\raggedright\arraybackslash}X>{\raggedright\arraybackslash}X@{}}
\toprule
\textbf{} & \textbf{word} & \textbf{same root} \\
\midrule
Kannada & \ka{ಬೆವರು} (\emph{bevaru}) & yes \\
Malayalam & \ml{വിയർപ്പ്} (\emph{viyarppŭ}) & yes \\
Hindi (neighbour) & \dv{पसीना} (\emph{pasīnā}) &  \\
English & sweat &  \\
\bottomrule
\end{tabularx}

\begin{grammarlens}[title={What you've built}]
A hospital, feeling sick, a headache, health, and sweat.
\end{grammarlens}

\subsection*{Your turn}
\begin{itemize}
\raggedright
  \item \emph{Say it:} \emph{viyarvai}
  \item \emph{Say it:} \emph{viyarvai}, once more
  \item \emph{Say it:} say \emph{viyarvai varukiṟatu}
  \item \emph{Recall:} say the Tamil for a headache, then the Tamil for health, then say \emph{viyarvai} again
\end{itemize}

\begin{tcolorbox}[breakable,colback=teal!4,colframe=teal!35!black,title={Before you move on}]
What does \ta{வியர்வை} mean? (\textquotedblleft{}Sweat\textquotedblright{}.) Say it once more.
\end{tcolorbox}
